\documentclass[%
 reprint,
 superscriptaddress, 
%groupedaddress,
%unsortedaddress,
%runinaddress,
%frontmatterverbose, 
%preprint,
%preprintnumbers,
%nofootinbib,
%nobibnotes,
%bibnotes,
 amsmath,amssymb,
 aps,
 pra,
%prb,
%rmp,
%prstab,
%prstper,
 floatfix,
]{revtex4-2}
% Rules of thumb that make it easier when you have to deal with journal submissions
%
% These are not all "good LaTeX" rules of thumb,
% but they help when working with antiquated journal submission systems.
%
% Have a single tex file
% Always have 0 errors and work hard to have 0 warnings - "a float is stuck" is probably the only warning it is ok to accept after an attempt or two to fix it
% Use only the automatic placement location for floats, figures, tables until the very last edit
% Have only alpha-numerical characters or underscores in your file names (no spaces).
% Do not use sub-directories, have everything in a single folder (do not put images in a separate folder).
% Use ~\cite and ~\ref (a fixed width non-breaking ~, not a generic whitespace in front of the command)
% Use \cite, not \citep or others
% Be consistent in abreviation, e.g. use only Fig. not Fig/Fig./fig/figure; only Eq. not (eq)/equation/eq./Eq
%
% Math:
% Avoid defining custom math commands
% In math formulas, use \mathrm for subscripts, unless the subscript is a dummy variable like an index
% In math formulas, use align, not eqarray, for systems of equations
% Always use numbered equations
%
% Images:
% Have a single image per figure (not two images included one after the other)
% Have your images be high DPI (>400) raster files or vector EPS files
% Avoid vector images if they are heatmaps or have more than hundred points or lines
% !IMPORTANT! Set the width of your images in your plotting software to the correct width in inches or cm not pixels, otherwise it would scale weirdly when imported in latex
% Set width=\linewidth for your images in your figures
% Set the font size of text in your images to one fixed value for all images, in pt, around 10-14pt is probably good
\usepackage{adjustbox}
\usepackage[table,dvipsnames]{xcolor}
\usepackage{array}
\usepackage{graphicx}% Include figure files
\usepackage{dcolumn}% align table columns on decimal point
\usepackage{bm}% bold math
\usepackage{enumitem}
\usepackage{multirow}
\usepackage{hyperref}
\hypersetup{hidelinks,hypertexnames=false}
\usepackage{amsthm}
\usepackage{amsmath}
\usepackage{tabularx}
\usepackage[linesnumbered,ruled,vlined]{algorithm2e}
\usepackage{stackengine}
\DeclareMathOperator{\kernel}{ker}
\DeclareMathOperator{\Hom}{Hom}
\newtheorem{definition}{Definition}
\newtheorem{theorem}{Theorem}
\newtheorem{lemma}{Lemma}
\usepackage{amsmath, amssymb}
\usepackage{hyperref}
\usepackage{amsmath,amsthm}
\usepackage{float}
\usepackage{epstopdf}
\usepackage{longtable}
\usepackage{mathtools}
\usepackage{colortbl}
\usepackage{booktabs}
\usepackage{mathtools,amssymb}
\usepackage{amsthm}
\usepackage[caption=false]{subfig}
% Unused nicematrix package removed for portable compilation.
\usepackage{tikz}
\usepackage{amsmath,amssymb}
\usetikzlibrary{positioning,arrows.meta,decorations.pathreplacing,calligraphy,calc}
\usetikzlibrary{cd}
\DeclareMathOperator{\im}{im}
\usepackage{tikz}
\usetikzlibrary{cd}
\usepackage{graphicx}% Include figure files
\usepackage{dcolumn}% align table columns on decimal point
\usepackage{bm}% bold math
\newtheorem{statement}[lemma]{Statement}
\usepackage{enumitem}
\usepackage{multirow}
\usepackage{hyperref}
\usepackage{float}
\usepackage{csquotes}% to use normal quotes
\usepackage[T1]{fontenc}
% allow long URLs (e.g. the repository URL in the bibliography) to break after hyphens
\makeatletter\AtBeginDocument{\g@addto@macro{\UrlBreaks}{\do\-}}\makeatother
\MakeOuterQuote{"}

\usepackage[linesnumbered, ruled, vlined]{algorithm2e}
\SetKw{KwDownTo}{downto}
\SetKwInput{KwInput}{Input}
\SetKwInput{KwOutput}{Output}
\usepackage{graphicx} 
\usepackage{mathtools}
\usepackage{amsmath}
\usepackage{amsfonts}
\usepackage{amssymb}
\usepackage{amsthm}
\usepackage{makecell}
\usepackage{physics}
% \newcommand\editcolor[1]{\textcolor{blue}{\textbf{#1}}}
\newcommand\editcolor[1]{{\color{blue}#1}}
\newcommand{\dhor}{\mathsf{h}}
\newcommand{\dver}{\mathsf{v}}
\newcommand{\hor}{\mathsf{v}}
\newcommand{\ver}{\mathsf{h}}
\newcommand{\lmp}[1]{\operatorname{L}_{#1}}
\newcommand{\ann}[1]{\operatorname{ann}(#1)}
\newcommand{\Code}{\mathcal{C}}
\newcommand{\ftwo}{\mathbb{F}_2}
\DeclareMathOperator{\id}{id}
\DeclareMathOperator{\Tor}{Tor}
\newtheorem{corollary}{Corollary}[theorem]
\newtheorem{proposition}{Proposition}[section]
\newtheorem{observation}[theorem]{Observation}
\newtheorem{remark}{Remark}
\theoremstyle{definition}
\newtheorem{example}{Example}
\theoremstyle{definition}
\def\id{{\mathchoice {\rm 1\mskip-4mu l} {\rm 1\mskip-4mu l} {\rm 1\mskip-4.5mu l} {\rm 1\mskip-5mu l}}}
\usepackage{tikz}
\usetikzlibrary{arrows.meta, positioning}
\usepackage{braket}

\usepackage{listings}
\usepackage{color}
\definecolor{mygreen}{rgb}{0,0.6,0}
\definecolor{myred}{rgb}{0.9,0.2,0.3}
\definecolor{mymauve}{rgb}{0.28,0,0.92}
\definecolor{mygray}{rgb}{0.5,0.5,0.5}

\lstdefinelanguage{JuliaREPL}{
  keywords={julia},
  keywordstyle=\color{mygreen}\bfseries,
  keywords=[2]{Stabilizer, GeneralizedStabilizer, naive_encoding_circuit, PauliOperator, Destabilizer, project, CliffordOperator, apply},
  keywordstyle=[2]\color{mymauve}\bfseries,
  identifierstyle=\color{black},
  sensitive=true,
  mathescape=true,
  comment=[l]{\#},
  morecomment=[s]{/*}{*/},
  commentstyle=\color{mygray}\ttfamily,
  stringstyle=\color{myred}\ttfamily,
  morestring=[b]',
  morestring=[b]"
}

\lstset{
   language=JuliaREPL,
   extendedchars=true,
   basicstyle=\footnotesize\ttfamily,
   showstringspaces=false,
   showspaces=false,
   tabsize=2,
   breaklines=true,
   showtabs=false
}

\usepackage{csquotes}% to use normal quotes

% a package for neat todo notes
% comment this out when finishing the manuscript
\usepackage[textsize=tiny]{todonotes}

% a setup to mark edited lines when resubmitting to a journal after review
% disable this when you want to hide edit colors
% colors from tableau 10 scheme
% if there are yet more rounds of reviews add yet another command
\definecolor{customcolorblue}{HTML}{4573ae}
%\newcommand\editcolor[1]{\textcolor{customcolorblue}{\textbf{#1}}}

% Uniform table typography. Adjust these settings once for every table.

\newcommand{\MMTableStyle}{%
  \color{black}%
  \fontsize{9}{11}\selectfont
  \renewcommand{\arraystretch}{1.2}%
  \setlength{\tabcolsep}{2pt}%
  \setlength{\extrarowheight}{2pt}%
  \setlength{\LTleft}{\fill}%
  \setlength{\LTright}{\fill}%
}


% Revision markup: replace the body by #1 to produce a clean manuscript.
\DeclareRobustCommand{\revisionblue}[1]{{\color{blue}#1}}

\begin{document}
\title{Generating algorithms for multivariate multicycle codes}
\author{F. A. Mian, O. Gwilliam, and S. Krastanov}
\begin{abstract}Pseudocode accompanying ``Multivariate Multicycle Codes for Complete Single-Shot Decoding'' (formerly Appendix A of the manuscript).\end{abstract}
\maketitle

In this appendix, we present Algorithm~\ref{alg:mm-css-quo} for constructing Multivariate Multicycle (MM) codes via Koszul complexes over quotient rings. The algorithm takes polynomials $F_1, \ldots, F_t$ in the quotient ring 
\begin{equation}
S = \mathbb{F}_2[x_1, \ldots, x_D] / \langle x_1^{\ell_1} - 1, \ldots, x_D^{\ell_D} - 1 \rangle,
\end{equation}
where $t \geq 2$ is the number of polynomials, {$D \geq 1$} is the number of variables, and $\ell_1, \ldots, \ell_D \in \mathbb{Z}_{\geq 1}$ are exponents defining the ideal $I = \langle x_1^{\ell_1} - 1, \ldots, x_D^{\ell_D} - 1 \rangle$. It then constructs the Koszul complex $K_\bullet([F_1, \ldots, F_t]; S)$ over $S$, where each boundary map $\partial_k$ is a matrix with entries in $S$. {Each entry is expanded in the regular representation of the finite abelian group. The resulting binary blocks are multilevel circulant matrices; for one cyclic factor they are ordinary circulant matrices.}

{Throughout the main text, $\partial_k$ denotes the matrix of the boundary map acting on column vectors, so that $P_X = \partial_q$, $P_Z = \partial_{q+1}^{\top}$, $M_X = \partial_{q-1}$ and $M_Z = \partial_{q+2}^{\top}$. \texttt{Oscar.jl} represents module homomorphisms by matrices acting on row vectors; the array $\partial[k]$ in Algorithms~\ref{alg:mm-css-quo} and~\ref{alg:mm-css} therefore stores the transpose $\partial_k^{\top}$, which is why the assignments there read $P_X \gets \partial[q]^{\top}$ and $P_Z \gets \partial[q+1]$. }Algorithm~\ref{alg:mm-css} presents an alternative ``symbolic'' approach: it constructs the Koszul complex over the polynomial ring $\mathbb{F}_2[y_1, \ldots, y_t]$ with standard boundary maps having entries in $\{0, y_1, \ldots, y_t\}$, then substitutes each $y_i$ with the circulant matrix corresponding to $F_i \in S$. Both methods employ Algorithm~\ref{alg:polynomial-to-circulant} to convert polynomials to circulant matrices. After conversion, we verify pairwise commutativity of the resulting matrices. The algorithms are implemented in the open-source \texttt{QuantumClifford.jl} library using the OSCAR computer algebra system [OSCAR] and the Julia programming language.
\begin{algorithm}[H]
    \DontPrintSemicolon
    \SetAlgoLined
    \KwInput{$\ell_1, \dots, \ell_D$, $F_1, \dots, F_t \in S = \mathbb{F}_2[x_1, \dots, x_D] / \langle x_i^{\ell_i} - 1, \cdots x_D^{\ell_D} -1\rangle$}
    \KwOutput{$P_X$, $P_Z$, $M_X$, $M_Z$}
    $s \gets \prod_{i=1}^D \ell_i, t \gets \text{length}(F_1, \dots, F_t)$\;
    \For{$i \gets 1$ \KwTo $t$}{
        $C_i \gets \text{PolynomialToCirculant}(F_i, \ell_1, \dots, \ell_D)$
    }
    \tcp{Check pairwise commutativity}
    \For{$i \gets 1$ \KwTo $t$}{
        \For{$j \gets i+1$ \KwTo $t$}{
            \textbf{assert} $C_i C_j \equiv C_j C_i \pmod{2}$
        }
    }
    \tcp{Construct Koszul complex over S}
    $K \gets \text{KoszulComplex}([F_1, \dots, F_t], S)$
    $\partial \gets \text{Array}(t)$\;
    \For{$k \gets 1$ \KwTo $t$}{
        $\partial_k \gets \text{map}(K, k)$
        $B_k \gets \text{matrix}(\partial_k)$
        $(r,c) \gets \text{size}(B_k)$\;
        $M_k \gets \mathbf{0}_{r s \times c s}(\mathbb{Z})$\;
        \For{$i \gets 1$ \KwTo $r$}{
            \For{$j \gets 1$ \KwTo $c$}{
                $f \gets B_k[i,j]$\;
                \If{$f \neq 0$}{
                    $C \gets \text{PolynomialToCirculant}(f, \ell_1, \dots, \ell_D)$\;
                    $M_k[(i-1)s+1:is, (j-1)s+1:js] \gets C$\;
                }
            }
        }
        $\partial[k] \gets M_k \bmod 2$\;
    }
    \For{$k \gets 2$ \KwTo $t$}{
        \textbf{assert} $\partial[k] \cdot \partial[k-1] \equiv 0 \pmod{2}$
    }
    \eIf{$t = 2$}{
        $P_X \gets \partial[1]^\top$, $P_Z \gets \partial[2]$\;
        $M_X \gets \varnothing$, $M_Z \gets \varnothing$\;
    }{
        $q \gets \lfloor t/2 \rfloor$\;
        $P_X \gets \partial[q]^\top$, $P_Z \gets \partial[q+1]$\;
        \If{$t \geq 4$}{$M_X \gets \partial[q-1]^\top$}
        \If{$t \geq 3$}{$M_Z \gets \partial[q+2]$}
    }    
    \textbf{assert} $P_X P_Z^\top \equiv 0 \pmod{2}$\;
    \If{$M_X \neq \varnothing$}{\textbf{assert} $M_X P_X \equiv 0 \pmod{2}$}
    \If{$M_Z \neq \varnothing$}{\textbf{assert} $M_Z P_Z \equiv 0 \pmod{2}$}
    \KwRet{$(P_X, P_Z, M_X, M_Z)$}
    \caption{MM code via Koszul complex over quotient ring}
    \label{alg:mm-css-quo}
\end{algorithm}

\begin{algorithm}
    \DontPrintSemicolon
    \SetAlgoLined
    \SetKwComment{Comment}{// }{}
    \KwInput{$\ell_1, \dots, \ell_D$, $F_1, \dots, F_t \in R = \mathbb{F}_2[x_1, \dots, x_D] / \langle x_i^{\ell_i} - 1, \cdots x_D^{\ell_D} -1\rangle$}
    \KwOutput{$P_X$, $P_Z$, $M_X$, $M_Z$}
    $s \gets \prod_{i=1}^D \ell_i, t \gets \text{length}(F_1, \dots, F_t)$\;
    \BlankLine
    \For{$i \gets 1$ \KwTo $t$}{
        $C_i \gets \text{PolynomialToCirculant}(F_i, \ell_1, \dots, \ell_D)$\;
    }
    \BlankLine
    \Comment*[l]{Check pairwise commutativity}
    \For{$i \gets 1$ \KwTo $t$}{
        \For{$j \gets i+1$ \KwTo $t$}{
            \textbf{assert} $C_i C_j \equiv C_j C_i \pmod{2}$
        }
    }
    \BlankLine
    \Comment*[l]{Construct Koszul complex boundary maps}
    $R, y \gets \mathbb{F}_2[y_1, \dots, y_t]$\;
    $K \gets \text{KoszulComplex}(y)$\;
    $\partial \gets \text{Array}(t)$\;
    \For{$k \gets 1$ \KwTo $t$}{
        $B_k \gets \text{matrix}(\text{map}(K, k))$ 
        $(r,c) \gets \text{size}(B_k)$\;
        $\partial_k \gets \mathbf{0}_{rs \times cs}(\mathbb{Z})$\;
        \For{$i \gets 1$ \KwTo $r$}{
            \For{$j \gets 1$ \KwTo $c$}{
                \If{$B_k[i,j] = y_m$}{
                    $\partial_k[(i-1)s+1:is, (j-1)s+1:js] \gets C_m$\;
                }
            }
        }
        $\partial[k] \gets \partial_k \bmod 2$\;
    }
    \For{$k \gets 2$ \KwTo $t$}{
        \textbf{assert} $\partial[k] \cdot \partial[k-1] \equiv 0 \pmod{2}$
    }
    \BlankLine
    \eIf{$t = 2$}{
        $P_X \gets \partial[1]^\top$, $P_Z \gets \partial[2]$\;
        $M_X \gets \varnothing$, $M_Z \gets \varnothing$\;
    }{
        $q \gets \lfloor t/2 \rfloor$\;
        $P_X \gets \partial[q]^\top$, $P_Z \gets \partial[q+1]$\;
        \If{$t \geq 4$}{$M_X \gets \partial[q-1]^\top$}
        \If{$t \geq 3$}{$M_Z \gets \partial[q+2]$}
    }    
    \textbf{assert} $P_X P_Z^\top \equiv 0 \pmod{2}$\;
    \If{$M_X \neq \varnothing$}{\textbf{assert} $M_X P_X \equiv 0 \pmod{2}$}
    \If{$M_Z \neq \varnothing$}{\textbf{assert} $M_Z P_Z \equiv 0 \pmod{2}$}
    \KwRet{$(P_X, P_Z, M_X, M_Z)$}
    \caption{MM code via Koszul complex over polynomial ring}
    \label{alg:mm-css}
\end{algorithm}

\begin{algorithm}
    \DontPrintSemicolon
    \SetAlgoLined
    \KwInput{Polynomial $F$, orders $\ell_1, \dots, \ell_D$}
    \KwOutput{{Multilevel circulant matrix (circulant for $D=1$)} $C \in \mathbb{F}_2^{s \times s}$ s.t. $s = \prod_{i=1}^D \ell_i$}
    \BlankLine
    $s \gets \prod_{i=1}^D \ell_i$\;
    $C \gets \mathbf{0}_{s \times s}(\mathbb{F}_2)$\;
    $f \gets \text{lift}(F)$\;
    \For{$j \gets 0$ \KwTo $s-1$}{
        $tmp \gets j$\;
        \For{$k \gets D$ \KwDownTo $1$}{
            $idxs[k] \gets tmp \bmod \ell_k$\;
            $tmp \gets \lfloor tmp / \ell_k \rfloor$\;
        }
        \For{$\mathit{term} \in \text{terms}(f)$}{
            $c \gets \text{coeff}(\mathit{term},1)$\;
            \If{$c = 0$}{continue}
            $e_k \gets \text{degree}(\mathit{term},k)$ for $k = 1, \dots, D$\;
            $row \gets 0$\;
            \For{$k \gets 1$ \KwTo $D$}{
                $row \gets row \cdot \ell_k + ((idxs[k] + e_k) \bmod \ell_k)$\;
            }
            $C[row+1,j+1] \gets C[row+1,j+1] + c$\;
        }
    }
    \KwRet{$C$}
    \caption{PolynomialToCirculant}
    \label{alg:polynomial-to-circulant}
\end{algorithm}

\end{document}
